\section{A decidable presentation of the real algebraic field}
\label{sec:scalar}
A scalar is a finite record. We construct the records, their primitive
recursive (PR) operations and exact comparisons. In the formulas below
rationals are pairs of a signed numerator and a positive denominator.
A polynomial is a coefficient list in ascending degree with trailing
zeros removed. Polynomial coefficients in a scalar record are rational.

\subsection{Reading a scalar}
For $p=\sum_i p_iX^i$ put
\[
 M(p,R)=\sum_i |p_i|R^i,\qquad D(p,R)=M(p'',R),\qquad R\ge1.
\]
A record $(p,l,r)$ has rational $l<r$. Set
$c=(l+r)/2$, $w=(r-l)/2$, $R=1+\max(|l|,|r|)$. Its validity test is
\begin{equation}
 p(l)p(r)<0,\qquad |p'(c)|>D(p,R)w. \label{eq:intervalcheck}
\end{equation}
In particular $(X-a,a-1,a+1)$ codes a rational $a$. Invalid input
can default to the record for zero. Define $p[u,v]$ by replacing
$X^i$ by $\sum_{j<i}u^{i-1-j}v^j$ for $i>0$, and replacing $1$ by zero.
Then $p(v)-p(u)=(v-u)p[u,v]$. For $|c|,|u|,|v|\le R$,
\begin{equation}
 |p[u,v]-p'(c)|\le D(p,R)\max(|u-c|,|v-c|).
 \label{eq:divdiff}
\end{equation}
Indeed, replacing the factors of a product one at a time costs at most
$(i-1)R^{i-2}$ times the maximum displacement, and the $i$th monomial
has $i$ such products. The constant and linear cases contribute zero.
Thus \eqref{eq:intervalcheck} gives a constant sign and a positive
rational lower bound for the absolute slope on the record interval.

Bisect, retaining the half with opposite endpoint signs, or freezing
at a midpoint with $p=0$. Start with $\lceil r-l\rceil+1$ steps and
use $n$ more for precision $n$. The nested closed intervals have width
at most $2^{-n}$. Write $x$ for their value. The assertion $p(x)=0$
has a finite justification at each precision. On an interval of width
$e$ with opposite endpoint signs, the absolute value at each endpoint is
at most $M(p',R)e$ by the divided-difference identity. Frozen intervals
give exact zero. Two roots in the same record interval agree by its
positive lower slope bound.

A \emph{provisional value} is a specified PR recipe for nested rational
intervals $I_n$ of width at most $2^{-n}$, possibly with numerical
parameters. Two such recipes agree if their intervals intersect at
every equal precision. This is transitive: disjoint outer intervals have
a positive rational gap, and the intermediate interval at finer precision
cannot meet both once its width is below that gap. Strict order has a pair of separated intervals
as witness. These are schematic predicates about numbers. For example,
for inputs bounded in absolute value by $R\ge1$, multiplying two
intervals of width $e$ gives width at most $2Re$. Shifting precision
by $\lceil 2R\rceil+1$ suffices for multiplication and by one for
addition. If an interval has rational distance $b>0$ from zero, intersect
later intervals with it and invert their endpoints. Widths increase by
at most $b^{-2}$. Rational interval formulas are inclusion-preserving,
so the shifted output intervals are nested. Polynomial identities hold
for the recipes: at any requested precision use rational midpoints and
propagate errors through the expressions with these bounds. If outputs
were disjoint, they would have a rational gap, contradicted by refinement
below that gap and the rational identity. For finite expressions of
variable size fold magnitude bounds forward and required precisions
backward through the operation list. There is no quantification over
programs or sequences in this use of recipes.

\subsection{Sign and conversion to a finite record}
Suppose a provisional $x$ has a known nonzero rational annihilator $P$.
Write $P=X^k G$ with $g=G(0)\ne0$, and put
\[
 a=\min\left(\tfrac14,\frac{|g|}{8(1+\sum_{i>0}|G_i|)}\right).
\]
When $|x|\le2a$, one has $|G(x)-g|\le2a\sum_{i>0}|G_i|\le|g|/4$.
Consequently $x^k=0$ there. If $k>0$ this implies $x=0$: an enclosure
excluding zero at distance $b>0$ would give $|x^k|\ge b^k>0$.
For $k=0$ the case is impossible. Read an enclosure of width $<a$.
If it strictly excludes zero its sign is the answer with a rational
separation witness. Otherwise $|x|\le a$ and the answer is zero.
Here and throughout $n=2+\lceil1/\epsilon\rceil$ suffices for
$2^{-n}<\epsilon$.

For conversion compute the monic squarefree part $h$ of $P$ and
polynomials $U,V$ with $Uh+Vh'=1$. We have $P\mid h^{\deg P}$ over
$\mathbb Q$, so $h(x)=0$ by the power argument. Here is one bounded
way to establish these polynomial facts. Factor the primitive integer
associate of $P$ by interpolation. For degree $d>0$ select $d+1$
nonzeros of $P(j)$ among $j=0,\ldots,2d$. There are enough by
successive exact division by linear factors. An integral divisor
has a signed integer divisor of $P(j)$ at each such point. Interpolate
each tuple of possible values and test integral coefficients and exact
divisibility. Rational factorizations are covered by Gauss's lemma:
if a prime divided all coefficients of a product of two primitive
integer polynomials, reduction modulo that prime would give a zero
product of nonzero polynomials over a field, contrary to their leading
terms. At most $d-1$ successful splits are required. Take the product
of the distinct monic factors. Each is coprime to its derivative by
the degree drop in characteristic zero. Extended Euclidean division
with at most $\deg h+1$ nonzero divisions gives the displayed
B\'ezout identity. One may instead compute $h$ by dividing $P$ by
$\gcd(P,P')$ and making the quotient monic. Nonzero constant $P$
cannot annihilate a value.

Choose an integer $R>3+\max_{u\in I_0}|u|$ and put
\[
 B=\max(1,M(V,R)),\qquad L=1+D(h,R).
\]
B\'ezout gives $|h'(x)|\ge1/B$. Take $\delta=2^{-s}<1/(32BL)$
and let $c$ be the midpoint of $I_s$. Then $|x-c|\le\delta$ and
$|h'(c)|\ge1/B-L\delta$. By \eqref{eq:divdiff} every divided
difference on $[c-2\delta,c+2\delta]$, including those with $x$,
has the same sign and absolute value $>1/(2B)$. The endpoints are on
opposite sides of $x$ at distance at least $\delta$. Their $h$ values
are strictly opposite in sign. The returned record
$(h,c-2\delta,c+2\delta)$ passes \eqref{eq:intervalcheck} because its
test radius is at most $R$. Uniqueness in its interval identifies its
reader with $x$. All requested precision was computed in advance from
the given polynomial and coarse enclosure. Different records can
represent one scalar.

\subsection{Operations and equality}
For records $x_1,\ldots,x_m$ with monic annihilators $p_i$ of positive
degrees $d_i$ form the finite rational algebra
\[
 E=\mathbb Q[T_1,\ldots,T_m]/(p_1(T_1),\ldots,p_m(T_m)).
\]
Its vectors use the basis $\prod_iT_i^{e_i}$, $0\le e_i<d_i$.
Independent monic divisions in the separate variables yield the unique
representatives. Multiplication by $v$ is a rational matrix of dimension
$\prod_i d_i$. Its monic characteristic polynomial annihilates $v$,
by applying Cayley--Hamilton to the vector of $1$. Here Cayley--Hamilton
is a coefficient identity from
$(YI-A)\operatorname{adj}(YI-A)=\det(YI-A)I$. Multiply the coefficient
identities on the left by the corresponding powers of $A$ and sum to
obtain cancellation. Determinants can use their finite permutation
sums. Substitution $T_i\mapsto x_i$ preserves reductions because each
$p_i(x_i)=0$, irrespective of relations among the different $x_i$.

Use $v=T_1+T_2,T_1-T_2,T_1T_2,-T_1$, respectively, for an annihilator
for each ring operation. Compute the provisional value with interval
arithmetic and convert as above. For a nonzero input use its separated
enclosure to invert, and the reversed polynomial $X^{\deg p}p(1/X)$
as annihilator. Inversion at zero is totalized by zero. Compare two
records by determining the sign of their difference. It is zero exactly
when the two readers agree: nested input pairs that intersect at every
precision have all interval differences containing zero, while a pair
separated at some precision remains separated under refinement. Thus
record equality is a PR setoid relation. The ordered-field laws hold
modulo this equality by the rational identities of the interval recipes.
In detail, trichotomy comes from the sign test, translation preserves
signed separation after a sufficiently accurate read of the common
summand, and positive lower bounds $b,c$ give the product lower bound
$bc$. Denote the resulting field by $F$.

\subsection{Polynomial IVT and exhaustive root lists}
For a nonconstant monic $f(Z)=Z^n+\sum_{i<n}x_i Z^i$ over $F$, take
the algebra $E$ just described and set
$\widetilde f(Z)=Z^n+\sum_{i<n}T_iZ^i$ (index the generators from zero).
Form the rational matrix of multiplication by $Z$ on
$E[Z]/(\widetilde f)$. Its monic characteristic polynomial $H(Y)$,
of degree $n\prod_i d_i$, satisfies
\begin{equation}
 f(y)=0\quad\Longrightarrow\quad H(y)=0
 \label{eq:elimroots}
\end{equation}
even for a provisional value $y$. Indeed $H(Z)=0$ in the finite
algebra, and evaluation preserves every defining relation when
$T_i\mapsto x_i$ and $Z\mapsto y$. Extra roots of coefficient
annihilators do not affect this implication.

For $s<t$ and $f(s)f(t)<0$, substitute $s+(t-s)X$, trim leading
zeros by the field comparison and make the result monic. Its degree
is positive. Bisect $[0,1]$. Each midpoint is rational and its polynomial
sign is computable in $F$. Freeze at an exact root. After $k$ steps the
width is at most $2^{-k}$. The provisional $y$ is a zero: the endpoint
sign argument for reading a record still applies, using a rational
upper bound on $M(f',1)$ for the transformed polynomial, obtained from
its coefficient enclosures. By \eqref{eq:elimroots} conversion returns
a record of $y$. Undo the substitution. The result is strictly interior
by endpoint nonvanishing and trichotomy. Positive scalars therefore
have positive square roots (apply this result to $X^2-b$ on $[0,b+1]$).
Odd-degree polynomials have zeros by leading-term dominance at
$\pm R$ for $R>1+\sum_{i<n}|f_i|$ after monic normalization. This
establishes the positive-square and odd-degree-root axioms of a real
closed field and the polynomial intermediate value property.

There is also a uniform PR list of all distinct roots in increasing order.
Return the empty list for a nonzero constant and flag the zero polynomial
as everywhere zero. Otherwise form $H$ by \eqref{eq:elimroots} after
monic normalization and take its monic squarefree part $h$. Choose
an integer $R>2+\sum_{i<\deg h}|h_i|$. All roots of $h$ are strictly
inside $[-R,R]$ by leading-term dominance. Compute $B,L$ as in the
conversion proof with radius $R+1$, and partition $[-R,R]$ by a dyadic
mesh $\delta<1/(32BL)$. Record every grid zero and every open cell
with opposite endpoint signs. Recover a provisional root in a crossing
cell by rational bisection and convert it. For any $F$ root $z$ of $h$,
locate it among the grid points by field comparisons. If it is in an
open cell, the B\'ezout bound at $z$ and \eqref{eq:divdiff} give constant
slope sign and absolute slope $>1/(2B)$ throughout that cell. Its endpoint
signs must be opposite, and the recovered root equals $z$. The same
estimate at a grid zero excludes another zero in either neighboring open
cell. This gives the complete ordered list for $h$. Keep only those
entries satisfying $f(z)=0$ in $F$. Every $F$ zero survives by
\eqref{eq:elimroots} and divisibility $H\mid h^{\deg H}$. Externally
the same grid argument over the ordinary reals identifies $F$ with all
real algebraic numbers. Every record has its real zero by bisection
and every real zero of a rational polynomial is found by the grid.

\subsection{Induction and uniformity}
Integers, rationals, coefficient lists, interval endpoints, permutations
and matrix tables are number codes. For example use $0$ for the empty
list and $\operatorname{cons}(a,b)=2^a(2b+1)$. Decoding, length and
indexed access are bounded by the list code. Builders enlarge codes by
primitive recursion. Euclidean division has degree fuel, determinants
and interpolation tuples range over the finite sets described above,
and the grid has $2R/\delta$ cells. A bisection calling a previously
defined PR field-sign operation is primitive recursion. Computing a field
sign uses rational interval reading and the scalar arithmetic, with no
calls to the algebraic-coefficient root finder. A failed bounded search
on malformed data returns a fixed default record or list. The proofs
above exclude that branch on valid arguments.

For the induction restriction, the signs in an initial portion of
a grid are a bounded Boolean fold. Polynomial division retains the
quotients, so its invariant is an identity of concrete coefficient
vectors. Cofactor identities are equalities of finite rational sums,
with cancellation paired by interchanging two indices. Width bounds
and endpoint signs in a bisection are rational inequalities or Boolean
field tests at the current stage. The product expansion in
\eqref{eq:divdiff} uses a finite telescoping sum. Factorization and
B\'ezout witnesses are coefficient vectors, and candidate divisors are
tested in bounded lists. For identities of provisional values, fix an
output precision before the calculation, propagate the rational bounds
through the finite operation list and apply induction to rational
data. For instance, in the power argument an enclosure allegedly
separated from zero supplies its own rational gap and hence a
precision that distinguishes its power from zero. Generalizing over
the requested precision is a logical inference after the finite
calculation. Induction over a property universally quantified over
precisions, and a universal evaluator for PR programs, are both
unnecessary.

For other constructions of PR field presentations see V.~Selivanov and
S.~Selivanova, \emph{Primitive recursive ordered fields and some
applications} (\url{https://arxiv.org/abs/2010.10189}). For a constructive
formalization see C.~Cohen, \emph{Construction of real algebraic numbers
in Coq} (\url{https://perso.crans.org/cohen/papers/realalg.pdf}).
